\section{Experiments}
\subsection{Multimodal Understanding}
\subsubsection{Evaluation Settings}
\paragraph{Benchmarks.} We evaluate LLaDA2.0-Uni across 21 multimodal understanding benchmarks, focusing on three core capabilities: general VQA, reasoning, and OCR/document understanding:
\begin{itemize}[leftmargin=16pt]
    \item \textbf{General Tasks:}
    MMStar~\citep{mmstar}, MMBench~\citep{Mmbench}, MME~\citep{fu2023mme}, HallusionBench~\citep{Hallusionbench},  RealWorldQA~\citep{RealworldQA} and SimpleVQA~\citep{Simplevqa}.
    \item \textbf{Reasoning Tasks:} 
    MMMU~\citep{yue2024mmmu}, and MMMU-Pro~\citep{yue2025mmmupro}, MathVista~\citep{mathvista}, We-Math~\citep{We-math}, MathVision~\citep{wang2024measuring}, and MathVerse~\citep{zhang2024mathverse}.
    \item \textbf{OCR\&Chart Tasks:} 
    ChartQA~\citep{Chartqa}, DocVQA~\citep{Docvqa}, InfoVQA~\citep{InfoVQA}, CharXiv~\citep{Charxiv}, OCRBench~\citep{Ocrbench}, and AI2D~\citep{AI2D}.
    \item \textbf{Other Multimodal Tasks:}
    CountBench~\citep{countbench}, VLRewardBench~\citep{VLRewardBench}, and V$^*$~\citep{vstar}.
\end{itemize}

\paragraph{Baselines.} 
To evaluate the multimodal performance of LLaDA2.0-Uni, we compare it against an extensive set of baselines.
%
We first compare LLaDA2.0-Uni with leading specialized VLMs, including Qwen2.5-VL-7B~\citep{bai2025qwen2} (AR-based) and LLaDA-V~\citep{you2025llada} (diffusion-based).
%
Furthermore, we evaluate LLaDA2.0-Uni against state-of-the-art unified models categorized into: 1) AR-based models, such as BAGEL~\citep{bagel} and InternVL-U~\citep{tian2026internvl}; and 2) diffusion-based models, such as Lumina-DiMOO~\citep{dimoo} and LLaDA-o~\citep{you2026lladao}.

\subsubsection{Multimodal Understanding Performance}
As shown in Table~\ref{tab:understanding}, LLaDA2.0-Uni demonstrates strong and comprehensive multimodal understanding capabilities.
%
Compared to existing diffusion-based unified models like Lumina-DiMOO and LLaDA-o, LLaDA2.0-Uni achieves significant improvements across all major categories, particularly in general VQA tasks (e.g., MMStar: 64.1 vs. 58.0) and complex reasoning tasks (e.g., MMMU: 50.1 vs. 44.9).
%
Furthermore, LLaDA2.0-Uni also delivers consistently high performance in challenging OCR and document understanding scenarios, where baselines like Lumina-DiMOO struggle significantly.
%
Most impressively, LLaDA2.0-Uni performs on par with state-of-the-art specialized VLMs such as Qwen2.5-VL-7B, even slightly outperforming it on specific metrics like MMStar (64.1 vs. 63.9) and CountBench (86.0 vs. 84.9).
%
Overall, these results confirm that LLaDA2.0-Uni closes the gap between unified diffusion architectures and top-tier specialized VLMs.

\input{tab/understanding}

\subsection{Text-to-Image Generation}
\subsubsection{Evaluation Settings}
\paragraph{Benchmarks.} 
We conduct a comprehensive evaluation using a suite of established public benchmarks, including GenEval~\citep{ghosh2023geneval}, DPG-Bench~\citep{hu2024ella}, One-IG Bench~\cite{chang2025oneig}, and UniGenBench~\cite{UniGenBench++} for general generative capabilities, as well as CVTG-2K~\citep{du2025textcrafter} for text rendering proficiency. 
%
To further evaluate reasoning-informed image generation, we also test LLaDA2.0-Uni on the WISE-Bench~\citep{niu2025wise}. 

\paragraph{Baselines.} To comprehensively assess the text-to-image generation capabilities of LLaDA2.0-Uni, we benchmark our model against a diverse spectrum of strong baselines.
%
We first compare LLaDA2.0-Uni with leading specialized generation models (Gen. Only), including diffusion models like FLUX.1 [Dev]~\citep{flux}, Lumina-Image 2.0~\citep{qin2025lumina}, Seedream 3.0~\citep{gao2025seedream}, Qwen-Image~\citep{wu2025qwen}, LongCat-Image~\citep{team2025longcat}, and Z-Image~\citep{cai2025z}, as well as AR-based models like Emu3~\citep{wang2024emu3} and Lumina-mGPT 2.0~\citep{xin2025lumina}.
%
Furthermore, we evaluate LLaDA2.0-Uni against state-of-the-art unified models categorized into: 
%
1) AR-based and hybrid (AR + Diff.) models, including Janus-Pro~\citep{Janus}, BAGEL~\citep{bagel}, and OmniGen2~\citep{wu2025omnigen2}, Hunyuan Image 3.0~\citep{cao2025hunyuanimage}, NextFlow~\citep{zhang2026nextflow}, and InternVL-U~\citep{tian2026internvl}; 
%
and 2) discrete diffusion (D-Diff.) and hybrid diffusion (D-Diff. + Diff.) models, including MMaDA~\citep{Mmada}, Lumina-DiMOO~\citep{dimoo} and LLaDA-o~\citep{you2026lladao}.

\input{tab/geneval}

\input{tab/dpg}

\begin{table}[!tb]
\centering
\renewcommand{\arraystretch}{1.2}
\setlength\tabcolsep{3pt}
\small
\newcolumntype{C}{>{\centering\arraybackslash}X}

\caption{Comparison of Text-to-Image Generation Ability on OneIG-EN Benchmark.}
\begin{tabularx}{\linewidth}{ccc| *{5}{C}|c} 
\toprule
\textbf{Type} & \textbf{Model} & \textbf{Arch.} & \textbf{Alignment} &\textbf{Text} &\textbf{Reasoning} &\textbf{Style} &\textbf{Diversity} &\textbf{~Overall$\uparrow$} \\
\midrule
\multirow{5}{*}{\rotatebox[origin=c]{90}{Gen. Only}} 
& FLUX.1 [Dev] & Diff. & 0.786 & 0.523 & 0.253 & 0.368 & 0.238 & 0.434 \\
& Lumina-Image 2.0 & Diff. & 0.819 &0.106 &0.270 & 0.354 &0.216 & 0.353 \\
& Seedream 3.0 & Diff. &0.818 &0.865 &0.275 &0.413 &0.277 &0.530 \\
& Qwen-Image & Diff. &0.882 &0.891 &0.306 &0.418 &0.197 &0.539 \\
& Z-Image-Turbo & Diff. &0.840 &0.994 &0.298 &0.368 &0.139 &0.528 \\

\midrule
\multirow{6}{*}{\rotatebox[origin=c]{90}{Unified}} 
& Janus-Pro & AR & 0.553 & 0.001 & 0.139 & 0.276 & 0.365 & 0.267 \\
& BAGEL & AR + Diff. & 0.769 & 0.244 & 0.173 & 0.367 & 0.251 & 0.361 \\
& OmniGen2 & AR + Diff. & 0.804 &0.680 &0.271 &0.377 &0.242 &0.475 \\
& InternVL-U & AR + Diff. &0.820 &0.740 &0.270 &0.400 &0.250 &0.500\\
% & MMaDA \\
& Lumina-DiMOO & Diff. &0.820& 0.550& 0.280& 0.400& 0.230 & 0.460  \\
% & LLaDA-o \\
\rowcolor{myblue}
&LLaDA2.0-Uni &D-Diff. + Diff. &0.882 & 0.661 &0.323 &0.400 &0.259 &0.505 \\
\bottomrule
\end{tabularx}
\label{table:oneig_en}
\vspace{-0.06in}
\end{table}

\begin{table}[!tb]
\centering
\renewcommand{\arraystretch}{1.2}
\setlength\tabcolsep{1pt}
\small
\newcolumntype{C}{>{\centering\arraybackslash}X}

\caption{Comparison of Text-to-Image Generation Ability on UniGenBench.}
\begin{tabularx}{\linewidth}{ccc|*{10}{C}|c} 
\toprule
\textbf{Type} & \textbf{Model} & \textbf{Arch.} & \textbf{Style} & \textbf{World} & \textbf{Attr.} & \textbf{Action} & \textbf{Relat.} & \textbf{Logic} & \textbf{Gram.} & \textbf{Comp.} & \textbf{Layout} & \textbf{Text} & \textbf{~Overall$\uparrow$} \\
\midrule
\multirow{5}{*}{\rotatebox[origin=c]{90}{Gen. Only}} 
&FLUX.1 [dev] & Diff. & 83.90 & 88.92 & 67.84 & 62.17 & 67.26 & 30.91 & 60.96 & 47.04 & 71.83 & 32.18 &61.30   \\

&Emu3-Gen & AR   & 86.80 & 77.06 & 51.39 & 40.11 & 49.75 & 19.32 & 52.94 & 36.86 & 44.78 &  1.15 &46.02 \\

&Seedream 3.0 &Diff. &98.19 &94.90 &84.62 &83.14 &80.18 &51.83 &60.30 &72.32 &88.74 &69.86 &78.41\\

&Qwen-Image & Diff. &94.70 &94.15 &87.93 &82.60 &80.08 &51.59 &60.96 &72.94 &86.57 &72.13 &78.36\\

&Z-Image & Diff. &96.80 &94.46 &82.48 &78.90 &80.20 &49.08 &68.98 &76.80 &84.89 &68.39 &78.10\\


\midrule
\multirow{6}{*}{\rotatebox[origin=c]{90}{Unified}} 
&Janus-Pro & AR & 90.80 & 86.71 & 67.74 & 64.26 & 68.40 & 37.05 & 64.44 & 62.11 & 72.01 &  2.59 &61.61   \\

&BAGEL &AR + Diff. & 90.20 & 85.60 & 67.74 & 61.98 & 70.69 & 30.23 & 66.44 & 58.12 & 76.49 &  7.76 &  61.53  \\

&OmniGen2 & AR + Diff. &91.90 &86.39 &72.12 &62.83 &68.27 &32.50 &59.89 &56.31 &71.64 &29.02 &63.09\\

&MMaDA & D-Diff. &82.40 &56.65 &48.39 &37.83 &50.25 &17.95 &55.75 &32.35 &30.22 &1.15 &41.35\\

&Lumina-DiMOO & D-Diff. &89.70 &90.03 &81.62 &71.12 &78.43 &45.45 &70.45 &73.32 &82.84 &25.57 &71.12\\

% & LLaDA-o \\
\rowcolor{myblue}
&LLaDA2.0-Uni & D-Diff. + Diff. &95.30  &93.67 &91.77 &85.65 &86.42 &63.99 &72.19 &85.82 &90.30 &31.23 &79.63\\
\bottomrule
\end{tabularx}
\label{table:unigen}
\vspace{-0.08in}
\end{table}


\begin{table}[!tb]
\centering
\renewcommand{\arraystretch}{1.2}
\setlength\tabcolsep{2pt}
\small
\newcolumntype{C}{>{\centering\arraybackslash}X}

\caption{Comparison of Text Rendering Ability on CVTG-2K Benchmark.}
\begin{tabularx}{\linewidth}{ccc| *{5}{C}| CC} 
\toprule
\multirow{2}{*}{\textbf{Type}} & \multirow{2}{*}{\textbf{Model}} &\multirow{2}{*}{\textbf{Arch.}} & \multicolumn{5}{c|}{\textbf{Word Accuracy~$\uparrow$}} & \multirow{2}{*}{\textbf{NED~$\uparrow$}} & \multirow{2}{*}{\textbf{CLIPScore~$\uparrow$}} \\
\cmidrule(lr){4-8}
& & & 2 regions & 3 regions & 4 regions & 5 regions & average & & \\
\midrule
\multirow{5}{*}{\rotatebox[origin=c]{90}{Gen. Only}} 
& FLUX.1 [Dev] & Diff. & 0.608 & 0.553 & 0.466 & 0.431 & 0.496 & 0.687 & 0.740 \\
& Seedream 3.0 & Diff. &0.628 &0.596 &0.604 &0.561 & 0.592 &0.853 &0.782\\
& Qwen-Image & Diff. & 0.837 &0.836 &0.831 &0.816 &0.829 &0.912 &0.802 \\
& LongCat-Image & Diff. &0.912 &0.873 &0.855 &0.831 &0.865 &0.936 &0.785 \\
& Z-Image-Turbo & Diff. &0.887 &0.866 &0.862 &0.834 &0.858 &0.928 &0.804\\
\midrule
\multirow{3}{*}{\rotatebox[origin=c]{90}{Unified}} 
& BAGEL & AR + Diff. &0.498 &0.391 &0.332 &0.291 &0.356 &0.657 &0.779 \\
& InternVL-U & AR + Diff. &0.729 &0.660 &0.618 &0.549 &0.623 &0.804 &0.816\\
& Lumina-DiMOO & D-Diff. & 0.723 &0.646 &0.571 &0.505 &0.590  &0.805 &0.831 \\
\rowcolor{myblue} 
& LLaDA2.0-Uni & D-Diff. + Diff. & 0.788  & 0.776 & 0.763 & 0.746  &0.765  & 0.911  & 0.818\\
\bottomrule
\end{tabularx}
\label{table:cvtg2k_rendering}
\vspace{-0.15in}
\end{table}

\begin{table}[!tb]
\centering
\renewcommand{\arraystretch}{1.2}
\setlength\tabcolsep{1pt}
\small
\newcolumntype{C}{>{\centering\arraybackslash}X}
\vspace{-0.1in}
\caption{Comparison of Reasoning-Informed Image Generation Ability on WISE Benchmark.}
\begin{tabularx}{\linewidth}{ccc| *{6}{C} |c} 
\toprule
\textbf{Type} & \textbf{Model} & \textbf{Arch.} & \textbf{Cultural} & \textbf{Time} & \textbf{Space} & \textbf{Biology} & \textbf{Physics} & \textbf{Chem.} & \textbf{~Overall$\uparrow$} \\
\midrule
\multirow{5}{*}{\rotatebox[origin=c]{90}{Gen. Only}} 
& SD3-Medium & Diff. &0.43 &0.50 &0.52 &0.41 &0.53 &0.33 &0.45 \\
& FLUX.1 [Dev] & Diff. & 0.48 & 0.58 & 0.62 & 0.42 & 0.51 & 0.35 & 0.50 \\
& Emu3-Gen & AR &0.34 &0.45 &0.48 &0.41 &0.45 &0.27 &0.39 \\
& Qwen-Image & Diff. &0.62 &0.63 &0.77 &0.57 &0.75 &0.40 &0.62 \\
& LongCat-Image &Diff. &0.66 &0.61 &0.72 &0.66 &0.72 &0.49 &0.65 \\

\midrule
\multirow{7}{*}{\rotatebox[origin=c]{90}{Unified}} 
& Janus-Pro & AR & 0.30 & 0.37 & 0.49 & 0.36 & 0.42 & 0.26 & 0.35 \\
& BAGEL & AR + Diff. & 0.44 & 0.55 & 0.68 & 0.44 & 0.60 & 0.39 & 0.52 \\
& NextFlow & AR + Diff. &0.62 &0.60 &0.70 &0.54 &0.58 &0.38 &0.59 \\
&InternVL-U & AR + Diff. &0.37 &0.51 &0.68 &0.39 &0.62 &0.39 &0.46 \\
& Lumina-DiMOO & D-Diff. & 0.35 & 0.43 &0.59 &0.31 &0.49 &0.34 &0.40 \\
% & LLaDA-o &D-Diff. + Diff.\\
\rowcolor{myblue}
&LLaDA2.0-Uni &D-Diff. + Diff. &0.54 &0.77 &0.82 &0.79 &0.87 &0.60 &0.68 \\
\rowcolor{myblue}
&~+ w/ thinking & D-Diff. + Diff. & 0.73  & 0.79 & 0.86 & 0.81 & 0.88 & 0.74 & 0.78 \\
\bottomrule
\end{tabularx}
\label{table:wise_benchmark}
\vspace{-0.06in}
\end{table}


\subsubsection{General Image Generation Performance}
\paragraph{GenEval.} Table~\ref{table:geneval} presents a comparison of model performance on the GenEval benchmark, which is designed to evaluate object-centric T2I generation using compositional prompts with diverse object attributes. 
%
LLaDA2.0-Uni demonstrates strong compositional capabilities, achieving an overall score of 0.89.
%
This performance is highly competitive, significantly outperforming all unified models and bridging the gap with top-tier generation-only models.
%
In particular, LLaDA2.0-Uni shows a clear advantage in spatial arrangement, securing the highest Position score (0.90) across all evaluated models.




\paragraph{DPG-Bench.} Table~\ref{table:dpg_benchmark} reports the text-to-image generation results on the DPG benchmark.
%
Notably, LLaDA2.0-Uni achieves the state-of-the-art overall score of 87.76 among unified models, outperforming strong baselines such as LLaDA-o (87.04) and HunyuanImage-3.0 (86.10).
%
Specifically, our model secures the highest scores on the Entity (93.55) and Other (94.04) sub-metrics.
%
Furthermore, LLaDA2.0-Uni delivers highly competitive performance even against specialized generation-only models, surpassing Z-Image-Turbo (84.86) and demonstrating robust text-image alignment capabilities.


\paragraph{One-IG Bench.} As shown in Table~\ref{table:oneig_en}, LLaDA2.0-Uni achieves a highly competitive overall score of 0.505 on OneIG-EN.
%
Notably, it outperforms all other unified models in Alignment (0.882) and Reasoning (0.323), reaching levels comparable to top dedicated generation models like Qwen-Image.
%
However, LLaDA2.0-Uni falls short of leading models in generating dense text, indicating an area for future improvement.

\paragraph{UniGenBench.} LLaDA2.0-Uni sets a new standard for unified models on UniGenBench (EN), achieving a top overall score of 79.63 (Table~\ref{table:unigen}).
%
It performs consistently well across all ten dimensions, showing a clear advantage in Logic (63.99) and Layout (90.30), where it even surpasses many specialized generation models.
%
These results demonstrate that LLaDA2.0-Uni effectively closes the performance gap between general-purpose unified models and top-tier specialized models.


\subsubsection{Text-Centric Image Generation Performance}
\paragraph{CVTG-2k.}
The CVTG-2K benchmark evaluates the capability of the model for text rendering across multiple regions.
%
As reported in Table~\ref{table:cvtg2k_rendering}, LLaDA2.0-Uni leads the unified models with an overall score of 0.765.
%
A key observation is its exceptional stability in multi-region text generation.
%
While baselines like BAGEL, Lumina-DiMOO, and InternVL-U show a sharp performance drop when the number of regions increases, our model experiences a much slower decline.


\subsubsection{Reasoning-Informed Image Generation Performance}
\paragraph{WISE-Bench.} 
The benefits of our large-scale multimodal pre-training are evident on the WISE benchmark (Table~\ref{table:wise_benchmark}), which evaluates complex semantic understanding and world knowledge in image generation.
%
LLaDA2.0-Uni achieves a strong overall score of 0.68, ranking first among all unified models and performing on par with generation-only models like LongCat-Image.
%
% It also matches the performance of generation-only models like Qwen-Image.
%
Notably, incorporating a reasoning mode yields an additional 10\% improvement. 
%
These results highlight the strong capability of LLaDA2.0-Uni in reasoning-informed generation.


\begin{table}[!tb]
\centering
\renewcommand{\arraystretch}{1.2}
\setlength\tabcolsep{1.5pt}
\footnotesize % 调小字号以适应多列宽度
\newcolumntype{C}{>{\centering\arraybackslash}X}
\caption{Comparison of Instruction-based Image Editing Ability on ImgEdit Benchmark.}
\begin{tabularx}{\linewidth}{cc| *{9}{C} |c} 
\toprule
\textbf{Type} & \textbf{Model} & \textbf{Add} & \textbf{Adjust} &\textbf{Extract} &\textbf{Replace} &\textbf{Remove} &\textbf{Back.} &\textbf{Style} &\textbf{Hybrid} &\textbf{Action}
&\textbf{Overall $\uparrow$} \\
\midrule
\multirow{4}{*}{\rotatebox[origin=c]{90}{Gen. Only}} 
& FLUX.1 Kontext & 4.25 & 4.15 & 2.35 & 4.56 & 3.57 & 4.26 & 4.57 & 3.68 & 4.63 & 4.00 \\
& Step1X-Edit &3.88 &3.14 &1.76 &3.40 &2.41 &3.16 &4.63 &2.64 &2.52 &3.06 \\
& Qwen-Image-Edit &4.32 &4.36 &4.04 &4.64 &4.52 &4.37 &4.84 &3.39 &4.71 &4.35 \\
& Z-Image-Edit &4.40 &4.14 &4.30 &4.57 &4.13 &4.14 &4.85 &3.63 &4.50 &4.30 \\

\midrule
\multirow{5}{*}{\rotatebox[origin=c]{90}{Unified}} 
& BAGEL & 3.56 & 3.31 & 1.70 & 3.30 & 2.62 & 3.24 & 4.49 & 2.38 & 4.17 & 3.20 \\
& OmniGen2 & 3.57 & 3.06 & 1.77 & 3.74 & 3.20 & 3.57 & 4.81 & 2.52 & 4.68 & 3.44 \\
& InternVL-U &4.13 &3.40 &2.27 &4.13 &3.39 &3.84 &4.77 &3.03 &4.05 &3.67\\
& Lumina-DiMOO &3.41 &2.38 &1.90 &3.26 &2.21 &2.11 &4.19 &2.26 &3.17 &2.77\\
\rowcolor{myblue} 
& LLaDA2.0-Uni &3.76 &4.16 &2.40 & 4.04 & 3.82 & 4.07 & 4.60 & 3.97 & 4.42 & 3.92 \\
\bottomrule
\end{tabularx}
\label{table:imgedit_results}
\vspace{-0.15in}
\end{table}

\begin{table}[!tb]
\centering
\renewcommand{\arraystretch}{1.2}
\setlength\tabcolsep{2pt}
\small
\newcolumntype{C}{>{\centering\arraybackslash}X}
\vspace{-0.1in}
\caption{Comparison of Instruction-based Image Editing Ability on GEdit Benchmark. Abbreviations: Semantic Consistency (G\_SC), Perceptual Quality (G\_PQ), and Overall Score (G\_O).}
\begin{tabularx}{\linewidth}{ccc| *{3}{C} | *{3}{C}} 
\toprule
\multirow{2}{*}{\textbf{Type}} & \multirow{2}{*}{\textbf{Model}} & \multirow{2}{*}{\textbf{Arch.}} & \multicolumn{3}{c|}{\textbf{GEdit-Bench-EN$\uparrow$}} & \multicolumn{3}{c}{\textbf{GEdit-Bench-CN$\uparrow$}} \\
\cmidrule(lr){4-6} \cmidrule(lr){7-9}
& & & \textbf{G\_SC} & \textbf{G\_PQ} & \textbf{G\_O} & \textbf{G\_SC} & \textbf{G\_PQ} & \textbf{G\_O} \\
\midrule
\multirow{5}{*}{\rotatebox[origin=c]{90}{Gen. Only}} 
& FLUX.1 Kontext & Diff. & 6.52 & 7.38 &  6.00 & - & - & - \\
& Step1X-Edit & Diff. &7.66 &7.35 &6.97 &7.20 &6.87 &6.86 \\
& Qwen-Image-Edit & Diff. &8.00 & 7.86 & 7.56 & 7.82 & 7.79 & 7.52\\
&LongCat-Image-Edit & Diff. &8.18 &8.00 &7.64 &8.08 &7.99 &7.60 \\
&Z-Image-Edit & Diff. &8.11 &7.72 &7.57 &8.03 &7.80 &7.54 \\
\midrule
\multirow{5}{*}{\rotatebox[origin=c]{90}{Unified}} 
& BAGEL & AR + Diff. & 7.36 & 6.83 & 6.52 & 7.34 & 6.85 & 6.50 \\
& OmniGen2 & AR + Diff. & 7.16 & 6.77 & 6.41 & - & - & - \\
& InternVL-U & AR + Diff. & - & - & 6.66 & - & - & -\\
& Lumina-DiMOO  & D-Diff. &- &- &3.91 &- &- & -\\
\rowcolor{myblue} 
& LLaDA2.0-Uni & D-Diff. + Diff. &6.68 &7.52 &6.61 &6.63 &7.67 &6.66 \\
\bottomrule
\end{tabularx}
\label{table:gedit_bench}
\vspace{-0.08in}
\end{table}

\begin{table}[!tb]
\centering
\renewcommand{\arraystretch}{1.2}
\setlength\tabcolsep{3pt}
\small
\newcolumntype{C}{>{\centering\arraybackslash}X}

\caption{Comparison of Multi-Reference Image Editing Ability on MICo-Bench.}
\begin{tabularx}{\linewidth}{cc| *{4}{C}|c} 
\toprule

\textbf{Model} & \textbf{Arch.} & \textbf{Object} &\textbf{Person} &\textbf{HOI} &\textbf{De\&Re} &\textbf{~Overall$\uparrow$} \\

\midrule
Qwen-Image-Edit & Diff. & 52.4  & 21.1 & 35.0 & 37.4 & 35.9 \\

BAGEL & AR + Diff. & 39.0 & 28.5 & 25.3 & 44.5 & 34.4 \\

OmniGen2 & AR + Diff. & 46.3 & 22.9 & 32.2 & 36.8 & 33.8 \\

Lumina-DiMOO & D-Diff. & 38.4 & 12.1 & 24.7 & 21.3 & 23.3 \\
\rowcolor{myblue}

LLaDA2.0-Uni &D-Diff. + Diff. &51.0 &32.8 &46.0 &54.4 &47.1  \\
\bottomrule
\end{tabularx}
\label{table:mico-bench}
\vspace{-0.15in}
\end{table}

\subsection{Image Editing}
\subsubsection{Evaluation Settings}
\paragraph{Benchmarks.}
We evaluate LLaDA2.0-Uni on two general instruction-based image editing benchmarks: ImgEdit-Bench~\citep{ye2025imgedit} and GEdit-Bench~\citep{liu2025step1x-edit}. 
%
Additionally, we provide qualitative comparisons against leading models on MICo-Bench~\citep{wei2025mico}, a challenging multi-reference image editing benchmark.

\paragraph{Baselines.} 
We evaluate LLaDA2.0-Uni against specialized editing models (FLUX.1 Kontext~\citep{labs2025flux1kontextflowmatching}, Step1X-Edit~\citep{liu2025step1x-edit}, Qwen-Image-Edit~\citep{wu2025qwen}, Z-Image-Edit~\citep{cai2025z}) and unified models (BAGEL~\citep{bagel}, OmniGen2~\citep{wu2025omnigen2}, InternVL-U~\citep{tian2026internvl}, Lumina-DiMOO~\citep{dimoo}).
%
For the multi-reference editing benchmark, we compare with BAGEL, Qwen-Image-Edit, and OmniGen2, as they are the only baselines that natively support multiple image inputs.




\subsubsection{General Image Editing Performance}
\paragraph{ImgEdit.} Table~\ref{table:imgedit_results} presents the instruction-based image editing performance on the ImgEdit benchmark.
%
Among unified models, LLaDA2.0-Uni achieves the best Overall score of 3.92, ranking first and significantly outperforming peers like OmniGen2 (3.44) and InternVL-U (3.67).
%
Notably, LLaDA2.0-Uni excels in the Adjust and Hybrid tasks, securing the highest scores within the unified category. 
%
These results underscore its robust capability to comprehend and execute intricate editing instructions.

\paragraph{GEdit-Bench.} 
To further assess complex image editing capabilities, we evaluate our model on the GEdit benchmark.
%
As detailed in Table~\ref{table:gedit_bench}, LLaDA2.0-Uni achieves solid overall scores across both English (6.61) and Chinese (6.66) evaluations.
%
A key highlight is its strong performance in the Perceptual Quality category, proving that the model can execute edits without sacrificing the visual quality of the original image.


\subsubsection{Multi-Reference Image Editing Performance}
\paragraph{MICo-Bench.} 
To evaluate the multi-reference image composition capabilities of LLaDA2.0-Uni, we conduct experiments on the MICo-Bench.
%
As shown in Table~\ref{table:mico-bench}, LLaDA2.0-Uni achieves the best overall performance, setting a new state-of-the-art on this benchmark with a score of 47.1.
%
It significantly outperforms strong baselines such as OmniGen2 (33.8) and Qwen-Image (35.9).
%
Notably, Lumina-DiMOO, which shares the same dLLM architecture, struggles significantly on this task, yielding an overall score of only 23.3.
%
This stark contrast clearly validates the effectiveness of our architecture and data pipeline.


\subsection{Interleaved}
\subsubsection{Interleaved Generation}
\paragraph{Benchmark Construction.} 
Existing models like Bagel and Nextflow support interleaved generation, yet a standard benchmark remains absent.
%
Datasets like ISG-BENCH~\citep{chen2024interleaved} and OpenING~\citep{zhou2025opening} contain interleaved cases but specialized tasks (e.g., 3D scene transformation) that most models cannot process.
%
To address this, we propose the InterGen benchmark.
%
As shown in Figure~\ref{fig:intergen}, it is structured into three main categories and various subcategories to comprehensively cover practical interleaved applications.
%
InterGen comprises 150 samples and utilizes advanced VLMs (Gemini-3 and Qwen3-VL) as a judge to assess performance across three dimensions: text coherence, text-image alignment, and ID consistency.


\begin{figure}[tb]
    \centering
    \includegraphics[width=1.0\linewidth]{figs/intergen.pdf}
    \caption{Overview of InterGen Benchmark.
    }
    \label{fig:intergen}
    \vspace{-0.1in}
\end{figure}

\begin{figure}[tb]
    \centering
    \includegraphics[width=1.0\linewidth]{figs/interleave_generation.pdf}
    \vspace{-0.25in}
    \caption{Qualitative Results on Interleaved Generation Task.}
    \label{fig:inter_generation}
    \vspace{-0.1in}
\end{figure}

\begin{table}[t]
\centering
\renewcommand{\arraystretch}{1.2}
\setlength\tabcolsep{2pt}
\small
\newcolumntype{C}{>{\centering\arraybackslash}X}

\caption{Comparison of Interleaved Generation Ability on InterGen Benchmark.}
\begin{tabularx}{\linewidth}{cc| *{2}{C} | *{2}{C} |  *{2}{C}} 
\toprule
\multirow{2}{*}{\textbf{Model}} & \multirow{2}{*}{\textbf{Arch.}} & \multicolumn{2}{c|}{\textbf{Story Telling}} & \multicolumn{2}{c|}{\textbf{Explanation}} & \multicolumn{2}{c}{\textbf{Event Forecasting}} \\
% \cmidrule(lr){4-6} \cmidrule(lr){7-9}
& & \textbf{Gemini$\uparrow$} & \textbf{Qwen3-VL$\uparrow$} & \textbf{Gemini$\uparrow$} & \textbf{Qwen3-VL$\uparrow$} & \textbf{Gemini$\uparrow$} & \textbf{Qwen3-VL$\uparrow$}\\
\midrule

Emu3.5 & AR + Diff.  &6.28 & 6.83 &6.19 &6.48 &5.08  & 5.75  \\
\rowcolor{myblue}
LLaDA2.0-Uni & D-Diff. + Diff. 	&6.42 &7.02 &6.22 &6.35 &5.19	&5.94\\
\bottomrule
\end{tabularx}
\label{table:inter_bench}
\end{table}

\paragraph{InterGen Benchmark.} We primarily compare our model with Emu3.5~\citep{cui2025emu35nativemultimodalmodels}, as other models capable of interleaved generation (such as NextFlow~\citep{zhang2026nextflow} and Mogao~\citep{liao2025mogao}) are not yet open-source.
%
As shown in Table~\ref{table:inter_bench}, LLaDA2.0-Uni generally outperforms Emu3.5 on the InterGen benchmark.
%
Specifically, it achieves higher scores in the Story Telling and Time Series Forecasting tasks, while demonstrating comparable performance in Explanation.
%
More visualizations of the interleaved generation are shown in Figure~\ref{fig:inter_generation}.

\begin{figure}[!t]
    \centering
    \vspace{-0.1in}
    \includegraphics[width=1.0\linewidth]{figs/interleave_reasoning.pdf}
    \vspace{-0.23in}
    \caption{Qualitative Results on Interleaved Reasoning Task.}
    \label{fig:inter_reason}
    \vspace{-0.13in}
\end{figure}

\subsubsection{Interleaved Reasoning}
Unified models have already achieved impressive results in standard image understanding and generation tasks.
%
Looking forward, exploring interleaved reasoning emerges as a critical bottleneck to overcome.
%
We believe this is an essential step toward building a strong synergy between visual generation and understanding.
%
In LLaDA 2.0-Uni, we conduct a preliminary exploration into interleaved reasoning capability.
%
As illustrated in Figure~\ref{fig:inter_reason}, through interleaved reasoning, our model successfully deduces logical strategies in in chess games and provides step-by-step solutions to physics problems.
%
These promising results give us great confidence to further expand interleaved reasoning capabilities in future research.


\subsection{Ablation Study}
\subsubsection{Analysis of SPRINT Acceleration} \label{sec:sprint_ablation}
\begin{table}[!b]
\centering
\vspace{-0.18in}
\renewcommand{\arraystretch}{1.22}
\caption{Performance and TPS comparison with and without SPRINT. (w/o SGLang)}
\label{tab:sprint}
\resizebox{\textwidth}{!}{%
\begin{tabular}{cccccccccccccc}
\toprule
\textbf{Method} & \textbf{Metric}
  & \textbf{AI2D} & \textbf{OCRB} & \textbf{MathVista}
  & \textbf{ChartQA} & \textbf{DocVQA} & \textbf{MMMU}
  & \textbf{MMStar} & \textbf{GenEval} & \textbf{DPG}
  & \textbf{Avg.} & \textbf{$\Delta$} \\
\midrule
\multirow{2}{*}{LLaDA2.0-Uni}
  & Score & 82.0 & 75.7 & 68.1 & 80.1 & 89.5 & 50.1 & 64.1 & 89.0 & 87.76 & 76.3 & -- \\
  & TPS   & 19.5 & 21.2 & 55.0 & 28.7 & 8.0  & 49.4 & 31.7 & 2.8  & 2.7  & 24.3 & -- \\
\midrule
\multirow{2}{*}{\ \ + SPRINT}
  & Score & 80.9 & 73.4 & 67.2 & 81.0 & 89.0 & 52.5 & 63.0 & 87.8 & 86.27 & 75.7 & $-$0.6 \\
  & TPS   & 42.9 & 36.0 & 75.0 & 62.3 & 27.6 & 52.2 & 49.2 & 5.1  & 7.8  & 39.8 & $\times$1.6 \\
\bottomrule
\end{tabular}%
}
\end{table}

Table~\ref{tab:sprint} evaluates SPRINT on nine multimodal benchmarks, covering both understanding (AI2D, OCRBench, MathVista, ChartQA, DocVQA, MMMU, MMStar) and generation (GenEval, DPG). 
%
SPRINT shifts the average score from 76.3 to 75.7 ($-$0.6) while accelerating generation from 24.3 to 39.8 TPS ($1.6\times$). 
%
Several benchmarks show score decreases; the most notable are OCRBench ($-2.3$) and DPG ($-1.5$).
%
OCRBench demands precise character-level prediction, where the lower threshold $\tau = 0.93$ may accept tokens before sufficient refinement. 
%
The speedup is largest on benchmarks with longer outputs, where the per-step savings from prefix pruning compound across many denoising iterations: DocVQA reaches $3.5\times$ (8.0 $\to$ 27.6), and ChartQA and AI2D both reach $2.2\times$.
%
SPRINT improves MMMU by +2.4 and ChartQA by +0.9. 
%
The non-uniform unmasking schedule concentrates refinement on uncertain positions, effectively increasing the denoising budget for difficult tokens without additional forward passes.
%
Furthermore, we are integrating LLaDA2.0-Uni with SGLang~\citep{Sglang} to accelerate inference, and the implementation will be made publicly available soon.


\subsubsection{Analysis of Diffusion Decoder} \label{sec:distilled_ablation}
Table~\ref{tab:distilled_perf} shows the performance and speed comparison between the Diffusion Decoder (50 steps) and Diffusion Decoder Turbo (8 steps) across five benchmarks. Speed is measured on a single GPU at $1024 \times 1024$ resolution with a batch size of 1, using BF16 precision.

\begin{table}[!htb]
\centering
\setlength\tabcolsep{10pt}
\renewcommand{\arraystretch}{1.0}
\caption{Performance and Speed Comparison between Diffusion Decoder and Diffusion Decoder Turbo.}
\resizebox{1.0\textwidth}{!}{%
\begin{tabular}{cc|ccccc}
\toprule
\textbf{Method} & \textbf{Speed (s / img)}
  & \textbf{GenEval} & \textbf{DPG} & \textbf{UniGenBench}
  & \textbf{OneIG-EN}
  & \textbf{WISE} 
  \\ 
\midrule
Diffusion Decoder (50 steps)
  & 32.95 & 0.89 & 87.76 & 79.63 & 0.505 & 0.68 \\
\midrule
Diffusion Decoder Turbo (8 steps)
  & 2.90 & 0.87 & 87.24 & 79.76 &0.500 & 0.68  \\
\bottomrule
\end{tabular}%
\label{tab:distilled_perf}
}
\end{table}

Through few-step acceleration, the Diffusion Decoder Turbo achieves an 11.4$\times$ speedup (from 32.95\,s/img to 2.90\,s/img) while maintaining competitive performance across all benchmarks.
%
Specifically, it retains a GenEval score of 0.87 (vs.\ 0.89 for the base) and a DPG score of 87.24 (vs.\ 87.76). 
%
Similarly, on the UniGenBench, OneIG-EN, and WISE benchmarks, Diffusion Decoder Turbo achieves performance on par with the original decoder.
%
Moreover, as shown in Figure~\ref{fig:decoder_comp}, the visual quality remains virtually indistinguishable.

\begin{figure}[!t]
    \centering
    \vspace{-0.1in}
    \includegraphics[width=1.0\linewidth]{figs/decoder_comp.pdf}
    \vspace{-0.23in}
    \caption{Visual Comparison of the Decoder and the Distilled Version.}
    \label{fig:decoder_comp}
    \vspace{-0.13in}
\end{figure}



